\section{Data preparation}

Data preparation is one of the most time-consuming phrase in comparison with others.
It comprises several activities for getting these data ready for later analytics phrases.
typically, a data preparation consists of:
\begin{itemize}
    \item \textbf{Collect data}: the whole data was collected for us so we mainly loads it into our workspace and use
    \item \textbf{Cleaning data}: publisher once claims the data was clean. But we still examines some of the importance criteria such as missing values, duplicated data and outliers.
    when it comes to imbalance data, we processed it later in the modeling sections.
    \item \textbf{Exploratory Data analysis}: this involved 
    summarizing the main characteristics of the data to gain better insights and 
    validate assumptions. We modified this section becomes the \textbf{Exploration and Visualization}
    so that we can emphrasize their characteristics and relationship between them later.
\end{itemize}


\subsection{Data Overview}

\subsubsection{Identification of Variable Types}

Variables were classified into numeric and categorical types based on their underlying measurement scales. Numeric variables were defined as continuous measures with meaningful quantitative interpretation, while categorical variables included nominal and ordinal attributes representing discrete or ordered categories.

\begin{table}[H]
\centering
\caption{Variable Types in the Dataset}
\begin{tabular}{p{6cm}|p{9cm}}
\hline
\textbf{Numeric Variables} & \textbf{Categorical Variables} \\
\hline
BMI, MentHlth, PhysHlth &
Diabetes\_012, HighBP, HighChol, CholCheck, Smoker, Stroke, HeartDiseaseorAttack, \\
& PhysActivity, Fruits, Veggies, HvyAlcoholConsump, AnyHealthcare, NoDocbcCost, \\
& GenHlth, DiffWalk, Sex, Age, Education, Income \\
\hline
\end{tabular}
\end{table}

Although the Age variable is numerically coded in the BRFSS dataset, it represents ordinal age groups rather than continuous age measured in years. Since the numeric codes do not correspond to equal intervals, treating Age as a continuous variable could lead to misleading statistical inference. Therefore, Age was treated as a categorical variable in subsequent analyses.

This classification strategy was adopted to avoid statistical bias and ensure valid interpretation in descriptive statistics and predictive modeling.

\subsubsection{Descriptive Statistics for Numeric Variables}

Descriptive statistics were used to summarize numeric variables and assess their distributional characteristics.

\begin{table}[H]
\centering
\caption{Descriptive statistics for numeric variables}
\label{tab:desc_numeric}

\resizebox{0.8\textwidth}{!}{%
\begin{tabular}{|l|r|r|r|r|r|r|}
\hline
Variable & Mean & SD & Median & IQR & Min & Max \\
\hline
BMI      & 28.382364 & 6.608694 & 27 & 7 & 12 & 98 \\
MentHlth & 3.184772  & 7.412847 & 0  & 2 & 0  & 30 \\
PhysHlth & 4.242081  & 8.717951 & 0  & 3 & 0  & 30 \\
\hline
\end{tabular}
}
\end{table}

The BMI variable exhibits a relatively high mean value along with a large standard deviation, indicating substantial variability in body weight within the study population. This reflects the coexistence of multiple weight groups, including a considerable proportion of individuals classified as overweight or obese.

In contrast, MentHlth and PhysHlth display asymmetric distributions with median values equal to zero and long right tails, indicating that most respondents did not experience days of poor health, while a small subgroup reported prolonged periods of poor physical or mental health. These distributional characteristics reflect heterogeneity in self-reported health status and suggest that these variables may play an important role in subsequent inferential analyses.

\subsubsection{Descriptive Statistics for Categorical Variables}

Descriptive statistics were used to summarize categorical variables and assess their distributional characteristics.

\begin{table}[H]
\centering
\caption{Descriptive statistics for categorical variables}
\label{tab:desc_categorical}

\footnotesize
\setlength{\tabcolsep}{5pt}
\renewcommand{\arraystretch}{1.1}

\resizebox{0.9\textwidth}{!}{%
\begin{tabular}{|l|r|r|r|r|}
\hline
\textbf{Variable} & \textbf{n\_levels} & \textbf{Mode} & \textbf{Mode Count} & \textbf{Mode (\%)} \\
\hline
Age                & 13 & 9 & 33244  & 13.10 \\
AnyHealthcare      & 2  & 1 & 241263 & 95.11 \\
CholCheck           & 2  & 1 & 244210 & 96.27 \\
Diabetes\_012        & 3  & 0 & 213703 & 84.24 \\
DiffWalk            & 2  & 0 & 211005 & 83.18 \\
Education           & 6  & 6 & 107325 & 42.31 \\
Fruits              & 2  & 1 & 160898 & 63.43 \\
GenHlth             & 5  & 2 & 89084  & 35.12 \\
HeartDiseaseorAttack& 2  & 0 & 229787 & 90.58 \\
HighBP              & 2  & 0 & 144851 & 57.10 \\
HighChol             & 2  & 0 & 146089 & 57.59 \\
HvyAlcoholConsump    & 2  & 0 & 239424 & 94.38 \\
Income               & 8  & 8 & 90385  & 35.63 \\
NoDocbcCost           & 2  & 0 & 233326 & 91.58 \\
PhysActivity          & 2  & 1 & 191920 & 75.65 \\
Sex                   & 2  & 0 & 141974 & 55.97 \\
Smoker                & 2  & 0 & 141257 & 55.68 \\
Stroke                & 2  & 0 & 243388 & 95.94 \\
Veggies               & 2  & 1 & 205841 & 81.14 \\
\hline
\end{tabular}
}
\end{table}

Most categorical variables exhibit class imbalance, particularly those related to health status and healthcare services. In contrast, demographic and socioeconomic variables show more evenly distributed class structures, reflecting greater heterogeneity across population groups. Therefore, class imbalance should be carefully considered in subsequent analyses and modeling tasks.

